\documentclass[11pt]{article}

\usepackage[utf8]{inputenc}
\usepackage[spanish,es-tabla]{babel}
\usepackage[a4paper,margin=2.0cm]{geometry}
\usepackage{booktabs}
\usepackage{tabularx}
\usepackage{array}
\usepackage{float}
\usepackage{xcolor}
\usepackage{hyperref}
\usepackage{amsmath}

\hypersetup{
  colorlinks=true,
  linkcolor=blue!55!black,
  urlcolor=blue!55!black
}

\newcommand{\img}[2]{%
  \begingroup
  \pdfximage width #1 {#2}%
  \pdfrefximage\pdflastximage
  \endgroup
}
\newcommand{\pathcell}[1]{\texttt{\footnotesize #1}}
\newcommand{\bad}{\textcolor{red!65!black}{no oculta}}
\newcommand{\explore}{\textcolor{blue!55!black}{exploratoria}}

\title{Extraccion de corridas: Chua no suave y carpeta arctan}
\author{Resumen desde JSON, CSV y figuras producidas}
\date{8 de mayo de 2026}

\begin{document}
\maketitle


La lectura principal es la siguiente: hay varias corridas donde si se ejecuto
la verificacion de ocultamiento, pero todas detectan trayectorias
\pathcell{TARGET}. Por tanto, ninguna de esas corridas queda soportada como
atractor oculto bajo la muestra usada. 

\section*{Corridas con verificacion de ocultamiento}

\begin{table}[H]
\centering
\caption{Corridas que si produjeron \pathcell{hidden\_target\_summary.json}.}
\scriptsize
\begin{tabularx}{\textwidth}{>{\raggedright\arraybackslash}p{2.1cm}cccclX}
\toprule
ID & Modelo & \(q\) & Muestras & TARGET & Contacto detectado & Conclusion \\
\midrule
V1 base &
piecewise & 0.999 & 84 & 5 &
\pathcell{E+:3, E-:2} &
\bad: muestra inicial ya conecta con el atractor desde \(E_+\) y \(E_-\). \\
V2 q=0.99 &
piecewise & 0.990 & 126 & 4 &
\pathcell{E+:1, E-:3} &
\bad: el candidato de \(q=0.99\) tambien tiene contacto desde equilibrios. \\
V3 q=0.99 ref. &
piecewise & 0.990 & 240 & 35 &
\pathcell{E+:17, E-:18} &
\bad: la refinacion fortalece la evidencia en contra del ocultamiento. \\
V4 E- ref. &
piecewise & 0.999 & 108 & 13 &
\pathcell{E-:13} &
\bad: la vecindad de \(E_-\) alcanza el candidato en varias muestras. \\
V5 E+ ref. &
piecewise & 0.999 & 108 & 40 &
\pathcell{E+:40} &
\bad: \(E_+\) es el contacto mas fuerte; no hay soporte para oculto. \\
V6 no suave &
piecewise & 0.990 & 72 & 2 &
\pathcell{E+:2} &
\bad: es el mejor candidato acotado, pero falla como oculto fuerte. \\
\bottomrule
\end{tabularx}
\end{table}

\noindent
 La corrida menos mala es
V6, con solo 2 contactos TARGET de 72 muestras, ambos desde \(E_+\).

\begin{figure}[H]
\centering
\img{0.78\textwidth}{../chua_no_suave_q0p99_run/chua_piecewise/final_pdf_figs/fig05b_hiddenness_overview.png}
\caption{No suave, \(q=0.99\): atractor candidato, equilibrios y trayectorias desde vecindades. Los contactos TARGET aparecen desde \(E_+\).}
\end{figure}

\begin{figure}[H]
\centering
\img{0.78\textwidth}{../chua_arctan/hidden_verify_q_0p99000_refined/hidden_verify/fig05b_hiddenness_q0p99000_refined_overview.png}
\caption{Verificacion refinada \(q=0.99\) encontrada en \pathcell{chua\_arctan}. El JSON la identifica como \pathcell{piecewise}; la refinacion detecta 35 contactos TARGET.}
\end{figure}

\begin{figure}[H]
\centering
\img{0.78\textwidth}{../chua_arctan/hidden_verify_refined/fig05b_hiddenness_refined_overview.png}
\caption{Resumen refinado \(q=0.999\) en \pathcell{chua\_arctan}: 53 contactos TARGET agregados entre \(E_+\) y \(E_-\).}
\end{figure}

\section*{Corridas exploratorias o sin buen resultado}

\begin{table}[H]
\centering
\caption{Exploraciones que no dan una validacion positiva de ocultamiento.}
\scriptsize
\begin{tabularx}{\textwidth}{>{\raggedright\arraybackslash}p{2.8cm}c>{\raggedright\arraybackslash}p{3.0cm}X}
\toprule
ID & Modelo & Evidencia producida & Por que no es buen resultado \\
\midrule
A1 arctan real &
arctan &
\pathcell{continuation summary} y figuras iniciales &
La continuacion diverge. El estado final en \(\epsilon=1\) es de orden
\(10^{10}\), y no se produjo \pathcell{hidden\_target\_summary.json}. \\
A2 sweep \(h=0.0075\) &
piecewise &
\begin{tabular}[t]{@{}l@{}}barrido \(q=0.99,0.995\)\\\(0.9975,0.999,1.0\)\end{tabular} &
Fue exploratorio: genera atractores candidatos por \(q\), pero no verifica
ocultamiento dentro del barrido. Los \(q=0.99\) y \(q=0.999\) luego fueron
verificados/refinados y resultaron \bad. \\
A3 sweep no suave &
piecewise &
\begin{tabular}[t]{@{}l@{}}barrido \(q=0.97,0.98,0.985\)\\\(0.99,0.995,1.0\)\end{tabular} &
Tambien fue exploratorio: las corridas terminaron \pathcell{ok}, pero no son
prueba de ocultamiento. El \(q=0.99\) completo fue verificado y dio 2 TARGET. \\
A4 completa no suave &
piecewise &
corrida completa con cuencas, FFT, bifurcaciones y verificacion &
Es el conjunto mas completo, pero la verificacion cuantitativa lo descarta como
oculto fuerte por los contactos desde \(E_+\). \\
\bottomrule
\end{tabularx}
\end{table}

\section*{Barridos \(q\)}

\begin{table}[H]
\centering
\caption{Barrido exploratorio en \pathcell{chua\_arctan/runs\_q\_sweep\_h0075\_lm8} (JSON: \pathcell{model=piecewise}).}
\scriptsize
\begin{tabular}{cccccc}
\toprule
\(q\) & Estado & \(a_0\) & \(x_f\) & \(y_f\) & \(z_f\) \\
\midrule
0.9900 & ok & 5.6218 & 4.4084 & -0.7921 & -7.0765 \\
0.9950 & ok & 5.7431 & 0.9894 & 1.3710 & 2.2759 \\
0.9975 & ok & 5.8006 & -2.9673 & 0.7480 & 6.7329 \\
0.9990 & ok & 5.8341 & 3.4246 & -0.9771 & -6.1702 \\
1.0000 & ok & 5.8561 & 0.8509 & 1.4704 & 1.2217 \\
\bottomrule
\end{tabular}
\end{table}

\begin{table}[H]
\centering
\caption{Barrido exploratorio en \pathcell{chua\_no\_suave\_q0p99\_run/q\_order\_sweep}.}
\scriptsize
\begin{tabular}{cccccc}
\toprule
\(q\) & Estado & \(a_0\) & \(x_f\) & \(y_f\) & \(z_f\) \\
\midrule
0.9700 & ok & 5.0114 & 2.5820 & -0.4106 & -6.5096 \\
0.9800 & ok & 5.3474 & 3.5202 & 1.4020 & -3.9536 \\
0.9850 & ok & 5.4906 & 3.4685 & -0.0721 & -7.0342 \\
0.9900 & ok & 5.6218 & -2.6717 & -1.5139 & 2.1532 \\
0.9950 & ok & 5.7431 & -3.4210 & 0.6799 & 7.0407 \\
1.0000 & ok & 5.8561 & -1.8416 & 0.3798 & 6.0088 \\
\bottomrule
\end{tabular}
\end{table}

\section*{Caso arctan}

Nyquist/DF encontro dos pares y selecciono la rama 0:
\[
\omega_0=2.0991692817,\qquad k=-1.3282768771,\qquad a_0=0.9019225739.
\]
Pero la continuacion no fue acotada. En \(\epsilon=1\) el estado final fue
\[
(-1.3415{\times}10^{10},\ -2.1666{\times}10^{10},\
-2.3332{\times}10^{10}),
\]
por lo que no hay verificacion oculta valida para ese modelo.

\begin{figure}[H]
\centering
\img{0.76\textwidth}{../chua_arctan_run/chua_arctan/final_pdf_figs/fig02_continuation_progress.png}
\caption{Modelo \pathcell{arctan}: progreso de continuacion. La rama seleccionada no entrega un candidato acotado.}
\end{figure}


\section*{Conclusion}


\begin{itemize}
  \item Un candidato no suave \(q=0.99\) relativamente cercano a ser interesante,
  pero rechazado por 2 contactos TARGET desde \(E_+\).
  \item Una familia piecewise refinada en \pathcell{chua\_arctan} con muchos mas
  contactos TARGET, especialmente en \(q=0.999\), por lo que queda peor como
  candidato oculto.
  \item Un arctan real en \pathcell{chua\_arctan\_run} que no llega a candidato
  valido porque diverge antes de la verificacion.
\end{itemize}

\end{document}
